\chapter{Sélection adverse et toxicité}\label{ch:hf:adverse-selection-and-toxicity}

Les exécutions d'un teneur de marché provenant d'une catégorie de traders perdent presque un tick par action dix secondes plus tard, et celles provenant des autres gagnent un demi-tick. Ses
cotations sont les mêmes pour les deux tant qu'il ne décide pas qu'elles ne devraient pas l'être. Dans le marché simulé de ce chapitre, 32~\% des exécutions d'un teneur de marché à lot unique proviennent
de traders informés, et ces exécutions ont un mark-out de $-0.95$ ticks par action contre $+0.55$ pour les autres. Quatre variables du carnet, visibles pendant que la
cotation reste affichée, prédisent le mark-out d'une exécution avec une corrélation de 0.18 hors échantillon (out-of-sample)~; retirer le côté que la prévision marque comme toxique fait baisser la part des exécutions informées
de 29~\% à 18~\%. La fréquence des retraits compte autant que le principe même de retirer~: au seuil naturel, le teneur de marché déplace sans cesse ses ordres
en queue de leur file d'attente.

\section{Mark-outs par contrepartie, plateforme et type d'ordre}

La sélection adverse (livre~1, chapitre~1) se mesure par les mark-outs (livre~2, chapitre~15~; livre~7, chapitre~23)~: le mouvement
d'un prix de référence après une exécution, signé contre le teneur de marché. Ce qui rend la mesure utile, c'est le regroupement. Un dealer qui répond à des demandes de
cotation (RFQ) connaît sa contrepartie et peut calculer le mark-out de chaque client (livre~9, chapitre~24, tarifie ainsi ses clients). Sur une bourse anonyme, le teneur
de marché ne sait pas qui l'a touché, mais il connaît la plateforme, le type d'ordre qui l'a touché (un ordre exécuter-ou-annuler routé entre plateformes, un ordre au
marché, un ordre qui a balayé plusieurs niveaux), la taille et l'heure de la journée. Battalio, Corwin et Jennings ont constaté, à partir de données propriétaires et publiques sur les actions américaines,
que les ordres à cours limité envoyés vers des plateformes versant de fortes rétrocessions ont été exécutés moins bien selon plusieurs mesures de qualité, dont la sélection adverse~: la
plateforme est un groupe dont le mark-out mérite d'être calculé.

Dans une simulation, la catégorie de la contrepartie est connue. Le simulateur enregistre, pour chaque exécution du teneur de marché, si l'agresseur était un
trader informé, et le teneur de marché ne le voit jamais~; la \cref{fig:hf:adverse-selection-and-toxicity:classes} répartit les exécutions d'un teneur de marché à un seul lot sur six
sessions de vingt minutes. Les exécutions informées se retournent contre le teneur en moins de deux secondes et poursuivent leur mouvement~: $-0.95$ ticks par action à dix secondes (erreur type
0.14). Les exécutions non informées conservent leur demi-fourchette~: $+0.55$ (0.10). Les traders informés représentent 32~\% des exécutions et fixent le signe du total, $+0.07$.

\begin{omfigure}
\begin{tikzpicture}
\begin{axis}[width=11cm,height=5.2cm,xlabel={\small secondes après l'exécution},ylabel={\small mark-out (ticks par action)},xmode=log,log basis x=10,
  xmin=0.4,xmax=70,xtick={0.5,1,2,5,10,20,60},xticklabels={0.5,1,2,5,10,20,60},ymin=-1.6,ymax=0.8,tick label style={font=\footnotesize},
  grid=major,grid style={gray!25},table/col sep=comma,legend style={font=\scriptsize,at={(0.5,-0.3)},anchor=north},legend columns=2]
\addplot[omThm,thick,mark=*,mark size=1.2pt] table[x=h,y=informed]{figdata/hft/06-adverse-selection-and-toxicity/classes.csv};
\addplot[omDef,thick,dashed,mark=square*,mark size=1.2pt] table[x=h,y=uninformed]{figdata/hft/06-adverse-selection-and-toxicity/classes.csv};
\draw[gray] (axis cs:0.4,0) -- (axis cs:70,0);
\legend{agresseur informé,agresseur non informé}
\end{axis}
\end{tikzpicture}

\omcaption{Mark-out des exécutions d'un teneur de marché à un seul lot par rapport au mid, selon la catégorie du trader qui les a touchées (la vérité du simulateur), sur six
sessions simulées de vingt minutes~; pondéré par les quantités. Données~: \texttt{fig\_toxicity.py} sur \texttt{hf\_toxicity.run}.}\label{fig:hf:adverse-selection-and-toxicity:classes}
\end{omfigure}

\section{Détection du flux informé}

Le teneur de marché ne peut pas voir la catégorie du prochain agresseur~; il peut en revanche observer les circonstances dans lesquelles arrivent les traders informés. La toxicité du flux
(livre~7, chapitre~9) désigne la part du flux qui inflige une sélection adverse aux fournisseurs de liquidité. Easley, López de Prado et O'Hara ont construit le
VPIN, une estimation de celle-ci à partir du déséquilibre des ordres en temps de volume, et l'ont proposé comme signal d'alerte de la volatilité induite par la toxicité. Andersen et Bondarenko
ont jugé le VPIN mauvais prédicteur de volatilité, culminant seulement après le flash crash, son contenu prédictif provenant d'une relation mécanique avec
l'intensité du trading, et ont mis en garde contre l'adoption de toute mesure de stress avant de la comparer à des indices de référence simples. La leçon pour un teneur de marché est de
prédire ce qui le concerne, à savoir le mark-out de sa propre prochaine exécution, et de juger le prédicteur à cette aune.

Dans un carnet à grand tick, les circonstances sont peu nombreuses. Le chapitre~5 a montré qu'un côté dont la file est petite par rapport à celle de l'autre va être consommé~; une
rafale de volume agressif contre un côté annonce la même chose~; un mid qui n'a pas bougé depuis un moment a plus de chances d'être périmé~; une fourchette plus large
que d'un tick signale que le carnet est en cours de recotation. Les variables du chapitre sont les quatre suivantes, calculées à partir du flux ordre par ordre pendant que la cotation reste affichée
(le \cref{lst:hf:adverse-selection-and-toxicity:features})~: le volume agressif contre le côté sur les deux dernières secondes (en lots), la part de l'autre
côté dans les deux meilleures files, le logarithme de un plus le nombre de secondes depuis le dernier changement du mid, et le fait que la fourchette dépasse un tick.

\section{Scores de toxicité}

\begin{definition}[Score de toxicité]\label{def:hf:adverse-selection-and-toxicity:score}
Un \emph{score de toxicité}\index{score de toxicité} est une estimation en temps réel, par le teneur de marché, de la sélection adverse qu'il subirait si un côté de ses
cotations était exécuté maintenant~: il s'agit ici de l'opposé d'une prévision du mark-out de l'exécution à un horizon fixé, à partir de variables observables pendant que la cotation reste affichée,
ajustée sur les exécutions passées du teneur de marché.
\end{definition}

Ajustée par moindres carrés sur les 919 exécutions de six sessions de calibration, la prévision du mark-out à dix secondes est
\[
\widehat{\mathrm{mo}}=1.43-0.029\,\text{flux}-1.79\,\text{part de la file du côté opposé}-0.072\ln(1+\text{calme})+0.028\,\text{large},
\]
en ticks par action. La part de l'autre côté dans les files porte presque tout le signal, comme le chapitre~5 l'avait prédit. La prévision est corrélée à 0.19 avec les
mark-outs réalisés sur lesquels elle a été ajustée et à 0.18 avec ceux de six autres sessions~: faible et stable. Triées en quintiles selon la prévision, les
exécutions des autres sessions réalisent des mark-outs allant de $-0.86$ à $+0.83$ ticks par action (la \cref{fig:hf:adverse-selection-and-toxicity:calibration})~; une corrélation de
0.18 sur des exécutions dont les mark-outs varient de plusieurs ticks suffit à les séparer.

\begin{omfigure}
\begin{tikzpicture}
\begin{axis}[width=11cm,height=5.2cm,xlabel={\small mark-out prévu, quintile des exécutions},ylabel={\small ticks par action},xtick={1,2,3,4,5},
  xmin=0.5,xmax=5.5,ymin=-1.3,ymax=1.3,tick label style={font=\footnotesize},grid=major,grid style={gray!25},table/col sep=comma,
  legend style={font=\scriptsize,at={(0.5,-0.3)},anchor=north},legend columns=2]
\addplot[omDef,thick,mark=*,mark size=1.3pt] table[x=group,y=pred]{figdata/hft/06-adverse-selection-and-toxicity/calibration.csv};
\addplot[omThm,thick,only marks,mark=square*,mark size=1.5pt,error bars/.cd,y dir=both,y explicit]
  table[x=group,y=real,y error=se]{figdata/hft/06-adverse-selection-and-toxicity/calibration.csv};
\draw[gray] (axis cs:0.5,0) -- (axis cs:5.5,0);
\legend{prévision moyenne,mark-out réalisé ($\pm$ une erreur type)}
\end{axis}
\end{tikzpicture}

\omcaption{Hors échantillon~: les exécutions du teneur de base sur six sessions de test, réparties en quintiles selon la prévision ajustée de leur mark-out à dix secondes~;
prévision moyenne et mark-out réalisé par quintile. Données~: \texttt{hf\_toxicity.calibration}.}\label{fig:hf:adverse-selection-and-toxicity:calibration}
\end{omfigure}

\section{Réponses~: élargir, retirer, décaler, refuser}

Un score ne vaut que par ce que le teneur de marché en fait. Il existe quatre réponses, et la structure de marché détermine celles qui sont disponibles.

\begin{definition}[Élargissement défensif]\label{def:hf:adverse-selection-and-toxicity:widening}
L'\emph{élargissement défensif}\index{elargissement défensif@élargissement défensif} (defensive widening) éloigne une cotation du \omterm{def:hf:fair-value:fair}{prix juste}, sur le côté ou sur les deux côtés où la sélection adverse est
attendue, en la laissant dans le carnet à un prix qui compense le mark-out attendu~; dans un carnet à un tick, cela revient à reculer d'un tick par rapport au
meilleur prix.
\end{definition}

Le fading retire entièrement la cotation~; du point de vue du routeur, le livre~10, chapitre~18 appelle la cotation qui disparaît à l'arrivée d'un ordre un
\emph{évanouissement de la cotation} (quote fade). Décaler les deux cotations dans le sens du mouvement attendu (chapitre~7). Refuser, c'est la version du dealer~: un dealer qui
connaît sa contrepartie et la tarife en conséquence, grâce à la segmentation des clients (livre~9, chapitre~24), ou la rejette sous last look (chapitre~9). Sur une
bourse anonyme, un teneur de marché à un tick dispose de deux de ces réponses~: retirer la cotation, ou élargir en reculant
(la \cref{fig:hf:adverse-selection-and-toxicity:responses}).

\begin{omfigure}
\begin{tikzpicture}[font=\small,x=1cm,y=0.5cm]
\foreach \dx/\name in {0/référence,4/retrait,8/élargissement} {
  \draw[gray!60] (\dx,0) -- (\dx+3,0);
  \node[font=\footnotesize,anchor=north] at (\dx+1.5,-1.3) {\name};
  \fill[omDef!35] (\dx+0.3,0) rectangle (\dx+0.9,3.0);
  \fill[omDef!35] (\dx+0.9,0) rectangle (\dx+1.4,1.4);
  \fill[omThm!35] (\dx+1.6,0) rectangle (\dx+2.1,4.0);
  \fill[omThm!35] (\dx+2.1,0) rectangle (\dx+2.7,2.2);
  \node[font=\scriptsize,anchor=north] at (\dx+0.6,-0.1) {99};
  \node[font=\scriptsize,anchor=north] at (\dx+1.15,-0.1) {100};
  \node[font=\scriptsize,anchor=north] at (\dx+1.85,-0.1) {101};
  \node[font=\scriptsize,anchor=north] at (\dx+2.4,-0.1) {102};
}
\draw[black,very thick] (1.15,1.4) -- (1.15,2.2);
\draw[black,very thick] (1.85,4.0) -- (1.85,4.8);
\draw[black,very thick] (9.85,4.0) -- (9.85,4.8);
\draw[black,very thick] (8.6,3.0) -- (8.6,3.8);
\draw[black,very thick] (5.85,4.0) -- (5.85,4.8);
\node[font=\scriptsize,align=center,anchor=west] at (11.3,2.6) {notre lot~:\\barre épaisse};
\end{tikzpicture}

\omcaption{Les deux réponses qu'un teneur de marché à un lot peut adopter dans un carnet à un tick lorsque son cours acheteur est jugé toxique (le carnet~: files acheteuses en bleu, files vendeuses
en rouge, prix en ticks~; notre lot sous forme de barre épaisse au-dessus d'une file). Base~: notre lot repose au meilleur cours acheteur, 100, et au meilleur cours vendeur, 101. Retrait~: le
cours acheteur est retiré. Élargissement~: le cours acheteur recule à 99, où il n'est presque jamais exécuté.}\label{fig:hf:adverse-selection-and-toxicity:responses}
\end{omfigure}

Le quoteur du \cref{lst:hf:adverse-selection-and-toxicity:quoter} applique l'une ou l'autre réponse chaque fois que le score d'un côté dépasse un seuil. Sur les six
sessions de test, avec un seuil égal à zéro (un côté est retiré chaque fois que son exécution est prévue comme perdante)~:

\begin{center}
\small
\setlength{\tabcolsep}{4pt}%
\begin{tabular}{@{}>{\raggedright\arraybackslash}p{2.6cm}>{\centering\arraybackslash}p{2.1cm}>{\centering\arraybackslash}p{2.1cm}>{\centering\arraybackslash}p{2.1cm}>{\centering\arraybackslash}p{2.1cm}>{\centering\arraybackslash}p{2.1cm}@{}}
\toprule
réponse & actions & part informée & mark-out à dix secondes & messages & P\&L (écart-type)\\
& par session & des exécutions & (ticks par action, erreur type) & par session & (\$ par session)\\
\midrule
aucune & 19\,167 & 29.4~\% & $-0.005$ (0.096) & 696 & $-75.00$ (178.23)\\
retrait & 6\,550 & 18.3~\% & 0.136 (0.199) & 1\,012 & 37.75 (134.30)\\
élargissement & 6\,683 & 20.0~\% & 0.112 (0.222) & 1\,513 & $-2.08$ (105.59)\\
\bottomrule
\end{tabular}
\end{center}

Les deux réponses éliminent un tiers de la part informée des exécutions, ce que le simulateur mesure exactement~; leurs mark-outs s'améliorent d'environ 0.1 à 0.15
ticks par action, mais les erreurs types sont de même ordre, et six sessions ne suffisent pas à classer le retrait par rapport à l'élargissement. Chacune des deux coûte les deux tiers du
volume et envoie beaucoup plus de messages. Faire varier le seuil montre pourquoi~:

\begin{center}
\small
\begin{tabular}{@{}lrrrr@{}}
\toprule
retrait au-dessus d'un score de & actions & part informée & mark-out (erreur type) & messages\\
\midrule
0 & 6\,550 & 18.3~\% & 0.136 (0.199) & 1\,012\\
0.2 & 10\,733 & 17.9~\% & 0.214 (0.146) & 952\\
0.4 & 15\,017 & 24.1~\% & 0.186 (0.123) & 841\\
0.6 & 18\,017 & 30.1~\% & $-0.060$ (0.103) & 740\\
\bottomrule
\end{tabular}
\end{center}

Avec un seuil de zéro, le score d'un côté franchit le seuil dans un sens puis dans l'autre, et chaque franchissement annule l'ordre et le renvoie en queue de
file lorsqu'il revient~; le chapitre~5 a montré que la queue d'une file est l'endroit où se trouvent les exécutions adverses. Un seuil de 0.2 maintient la part informée
à un niveau bas, avec 64~\% de volume en plus. À partir de 0.4, le teneur de marché se retire trop rarement et la part informée remonte.

\section{La toxicité n'est pas permanente}

Les exécutions du même quoteur dans la fenêtre d'annonce de chaque session de calibration, ainsi qu'avant et après cette fenêtre~:

\begin{center}
\small
\begin{tabular}{@{}lccc@{}}
\toprule
& avant l'annonce & pendant l'annonce & après l'annonce\\
\midrule
exécutions par minute & 6.97 & 17.89 & 6.69\\
part informée des exécutions & 34.7~\% & 23.6~\% & 32.3~\%\\
mark-out à dix secondes (ticks par action) & 0.025 & 0.059 & 0.132\\
\bottomrule
\end{tabular}
\end{center}

Sur ce marché, la fenêtre d'annonce apporte trois fois plus de flux de bruit, ainsi qu'un prix efficient qui s'ajuste plus vite~; les exécutions pendant cette fenêtre sont donc moins toxiques, et non
plus~: une règle qui élargit à chaque fenêtre d'annonce abandonnerait le meilleur flux de la séance. Le point se généralise~: la toxicité est une propriété du flux
à un instant donné, et non d'un instrument, d'une plateforme ou d'un client pour toujours. Un score est recalibré sur les exécutions récentes, sa décroissance se mesure comme pour tout prédicteur
(livre~7, chapitre~13), et une contrepartie classée toxique est réexaminée.

\section{Tutoriel~: calculer le score de la prochaine exécution}\label{tut:hf:adverse-selection-and-toxicity:tut}

\begin{tutorial}
\textbf{Objectif.} Mesurer qui se trouve de l'autre côté des exécutions d'un teneur de marché, construire un \omterm{def:hf:adverse-selection-and-toxicity:score}{score de toxicité} à partir de variables observables, et tester le retrait et
l'élargissement sur ce score. \textbf{État final~:} les trois tableaux et les \cref{fig:hf:adverse-selection-and-toxicity:classes,fig:hf:adverse-selection-and-toxicity:calibration}.
\begin{tutsteps}
\item \textbf{Variables issues du flux.} \texttt{firm.toxicity.FlowTracker} se met à jour à chaque message (le harnais transmet désormais le message à l'origine de chaque
mise à jour, \texttt{ctx.last}).
\omcode{code/firm/toxicity/firm_toxicity.py}{58}{75}{Flux signé pour chaque côté, part de la file du côté opposé, temps d'inactivité et fourchette.}\label{lst:hf:adverse-selection-and-toxicity:features}
\item \textbf{Calibrer.} \texttt{hf\_toxicity.fitted()} régresse les mark-outs à dix secondes des exécutions du quoteur de base sur les variables à chaque
exécution, sur six sessions de calibration.
\item \textbf{Réagir.} \texttt{ToxicQuoter} retire ou élargit un côté dont le score dépasse le seuil.
\omcode{code/firm/toxicity/firm_toxicity.py}{115}{137}{Calculer le score de chaque côté à chaque mise à jour~; retirer ou reculer le côté toxique.}\label{lst:hf:adverse-selection-and-toxicity:quoter}
\item \textbf{Évaluer} avec \texttt{responses()} et \texttt{fade\_sweep()} sur six autres sessions~; \texttt{by\_class()} et \texttt{by\_period()} utilisent
la vérité terrain du simulateur.
\end{tutsteps}
\textbf{Pour aller plus loin.} Ajouter une hystérésis (retrait au-dessus de 0.2, retour sous 0.1) et mesurer les messages économisés~; calculer la toxicité de tout le carnet
avec le VPIN et la comparer au score au niveau de chaque exécution.
\end{tutorial}

\section{Construction~: la boîte à outils de la toxicité}\label{bld:hf:adverse-selection-and-toxicity:toxicity}

\begin{build}
\bfield{Rôle} Mark-outs regroupés, variables de toxicité en flux continu, un score calibré et des réponses que tout quoteur peut envelopper.
\bfield{Interface} \texttt{grouped(markouts, qty, groups)}~; \texttt{FlowTracker(window).update(t, last, ext)}, \texttt{.features(side)}~;
\texttt{ToxicityModel.fit(X, y)}, \texttt{.predict}, \texttt{.score}~; \texttt{ToxicQuoter(model, mode, threshold, size, limit, record)}. Dans
\texttt{firm.mmharness}~: \texttt{ctx.last} (le message à l'origine de chaque mise à jour) et l'entrée \texttt{informed} de \texttt{Result.extra} (évaluation uniquement).
\bfield{Règles} Les variables n'utilisent que ce que le flux a montré à l'instant considéré. Un score est calibré sur des sessions différentes de celles sur lesquelles il est évalué. Les
classes de contreparties du simulateur n'atteignent jamais le quoteur.
\bfield{Tests d'acceptation} \texttt{code/firm/toxicity/tests/}~: variables calculées à la main sur un flux de trois messages (fenêtre, temps d'inactivité, fourchette)~; la calibration
retrouve un mark-out linéaire injecté~; moyennes regroupées à la main. \texttt{code/firm/mmharness/tests/}~: le quoteur voit les messages d'ajout, d'annulation et d'exécution~;
le journal des contreparties comporte une entrée par exécution et les deux classes.
\bfield{Pour aller plus loin} Hystérésis~; un score logistique de la probabilité informée~; des scores par plateforme lorsque \texttt{firm.exchsim} fait tourner plusieurs plateformes.
\end{build}

\begin{omsources}
\item D.~Easley, M.~López de Prado et M.~O'Hara, «~Flow toxicity and liquidity in a high-frequency world~», \emph{Review of Financial Studies}
25(5), 2012.
\item T.~G.~Andersen et O.~Bondarenko, «~VPIN and the flash crash~», \emph{Journal of Financial Markets} 17, 2014.
\item R.~Battalio, S.~A.~Corwin et R.~Jennings, «~Can brokers have it all~? On the relation between make-take fees and limit order execution
quality~», \emph{Journal of Finance} 71(5), 2016.
\end{omsources}

\section{Exercices}

\begin{exercise}[$\star$]\label{exo:hf:adverse-selection-and-toxicity:1}
32~\% des exécutions d'un teneur de marché sont informées, avec un mark-out de $-0.95$ tick~; les autres ont un mark-out de $+0.55$. Quel est le mark-out de l'ensemble de ses exécutions~?
\end{exercise}

\begin{exercise}[$\star$]\label{exo:hf:adverse-selection-and-toxicity:2}
Avec la prévision ajustée, calculer le score du cours acheteur lorsque 2 lots ont été vendus contre lui au cours des deux dernières secondes, et que le cours vendeur détient 80~\% des deux meilleures
files, que le mid a bougé pour la dernière fois il y a une demi-seconde et que la fourchette vaut un tick. Le teneur de marché déclenche-t-il un évanouissement de la cotation (quote fade) à un seuil de 0~? Et à 0.2~?
\end{exercise}

\begin{exercise}[$\star$]\label{exo:hf:adverse-selection-and-toxicity:3}
Une règle rapporte 0.2 tick par action sur 10\,000 actions par séance, avec un tick d'un centime. Quel est son edge sur dix secondes par séance~?
\end{exercise}

\begin{exercise}[$\star\star$]\label{exo:hf:adverse-selection-and-toxicity:4}
Pourquoi l'évanouissement et l'élargissement envoient-ils tous deux plus de messages que le teneur de marché de base~?
\end{exercise}

\begin{exercise}[$\star\star$]\label{exo:hf:adverse-selection-and-toxicity:5}
Pourquoi la part de l'autre côté dans les files porte-t-elle presque toute la prévision sur ce marché~?
\end{exercise}

\begin{exercise}[$\star\star$]\label{exo:hf:adverse-selection-and-toxicity:6}
D'après la \cref{fig:hf:adverse-selection-and-toxicity:calibration}, quel est le mark-out réalisé du cinquième des exécutions associé à la meilleure prévision, et
la différence avec le cinquième le plus faible est-elle supérieure à ses erreurs types~?
\end{exercise}

\begin{exercise}[$\star\star\star$]\label{exo:hf:adverse-selection-and-toxicity:7}
\emph{Programmation.} Lancer le teneur de marché avec évanouissement et hystérésis, \texttt{hysteresis(0.2, 0.0)}~: un côté est retiré de la cotation lorsque son score dépasse 0.2, et il n'est rétabli que
lorsque son score passe sous 0. Comparer ses messages, son volume et sa part d'exécutions informées avec un évanouissement au-dessus de 0.2 sans hystérésis.
\end{exercise}

\begin{exercise}[$\star\star\star$]\label{exo:hf:adverse-selection-and-toxicity:8}
\emph{Trouver la faille.} «~Le VPIN est à son plus haut de l'année~: élargir chaque cotation sur chaque plateforme jusqu'à ce qu'il baisse.~»
\end{exercise}

\section{Problème~: Qui est de l'autre côté}

\begin{problem}[{Problème du week-end --- qui est de l'autre côté}]\label{pb:hf:adverse-selection-and-toxicity:1}
Un teneur de marché à un lot, sur une action à un tick, veut cesser de négocier avec les personnes qui en savent plus que lui, sans pour autant arrêter de négocier.

\medskip\noindent\textbf{Partie I --- Mesurer.}
\begin{enumerate}
\item Selon quels groupes un teneur de marché peut-il ventiler son mark-out sur une bourse anonyme, et sur une plateforme de demande de cotation (RFQ)~?
\item Qu'ont constaté Battalio, Corwin et Jennings au sujet des plateformes~?
\item Donner les mark-out des exécutions informées et des exécutions non informées, ainsi que la part des exécutions informées.
\item Pourquoi le mark-out total est-il si proche de zéro~?
\end{enumerate}

\medskip\noindent\textbf{Partie II --- Prédire.}
\begin{enumerate}[resume]
\item Qu'est-ce que le VPIN, et qu'ont constaté Andersen et Bondarenko~?
\item Définir un \omterm{def:hf:adverse-selection-and-toxicity:score}{score de toxicité} et énumérer les variables du chapitre.
\item Donner la prévision ajustée ainsi que ses corrélations dans l'échantillon (in-sample) et hors échantillon (out-of-sample).
\item Que montre la calibration par cinquièmes~?
\end{enumerate}

\medskip\noindent\textbf{Partie III --- Réagir.}
\begin{enumerate}[resume]
\item Définir l'\omterm{def:hf:adverse-selection-and-toxicity:widening}{élargissement défensif} et le distinguer de l'évanouissement, du décalage des cotations et du refus de coter.
\item Quelles réponses un teneur de marché à un tick a-t-il sur une bourse anonyme~?
\item Donner la part de chacune des trois réponses, leur part d'exécutions informées et leurs mark-out.
\item Pourquoi six séances ne permettent-elles pas de classer l'évanouissement par rapport à l'élargissement~?
\end{enumerate}

\medskip\noindent\textbf{Partie IV --- Le verdict.}
\begin{enumerate}[resume]
\item Énoncer le \emph{résultat nommé}~: la variation de la part des exécutions informées et du mark-out à dix secondes lorsque le teneur de marché déclenche un évanouissement de ses cotations selon son score, à
un seuil de 0.2, ainsi que le volume abandonné.
\item Pourquoi un seuil de zéro est-il moins bon qu'un seuil de 0.2~?
\item Que se passe-t-il au-delà de 0.4~?
\item Que montre la fenêtre autour des annonces sur la toxicité au fil du temps~?
\item Comment re-tester une contrepartie classée comme toxique~?
\item Comment les réponses changeraient-elles sur un actif à petit tick~?
\item Que ferait un dealer avec le même score~?
\item En une phrase~: à quoi sert un \omterm{def:hf:adverse-selection-and-toxicity:score}{score de toxicité}~?
\end{enumerate}
\end{problem}

\section{Questions d'entretien}

\begin{interviewq}[$\star$ \iqroles{trader}]\label{iq:hf:adverse-selection-and-toxicity:1}
Vos exécutions sur une plateforme ont un mark-out nettement moins bon que sur une autre. Quelles sont trois explications possibles, et comment les distingueriez-vous~?
\end{interviewq}

\begin{interviewq}[$\star\star$ \iqroles{chercheur}]\label{iq:hf:adverse-selection-and-toxicity:2}
Vous n'observez que les mark-out des exécutions que vous avez obtenues. Comment cela biaise-t-il un modèle de toxicité, et que pouvez-vous faire pour y remédier~?
\end{interviewq}

\begin{interviewq}[$\star\star$ \iqroles{chercheur}]\label{iq:hf:adverse-selection-and-toxicity:3}
Une prévision de toxicité est corrélée à 0.18 avec les mark-out réalisés. Est-elle utile~? Que devez-vous savoir d'autre~?
\end{interviewq}

\begin{interviewq}[$\star\star$ \iqroles{développeur}]\label{iq:hf:adverse-selection-and-toxicity:4}
Un score franchit son seuil plusieurs fois par seconde. Quel effet cela a-t-il sur le teneur de marché, et comment le corrigez-vous dans le code~?
\end{interviewq}

\begin{interviewq}[$\star\star$ \iqroles{risque}]\label{iq:hf:adverse-selection-and-toxicity:5}
Un teneur de marché soumis à des obligations de cotation veut déclencher un évanouissement de ses cotations chaque fois que son score est élevé. Qu'est-ce qui le contraint~?
\end{interviewq}

\begin{interviewq}[$\star\star\star$ \iqroles{chercheur}]\label{iq:hf:adverse-selection-and-toxicity:6}
Une fraction $\pi$ des exécutions sont informées et perdent $L$~; les autres gagnent $G$. Un score signale les exécutions informées avec une probabilité $a$ et les exécutions non informées avec
une probabilité $b$. En déclenchant un évanouissement pour chaque exécution signalée, à quelle condition l'espérance du mark-out par exécution s'améliore-t-elle, et que devient le volume~?
\end{interviewq}
